\begin{minipage}[b]{0.4\linewidth}
    \centering
    \begin{tikzpicture}[->,>=stealth',node distance=5cm]
      \tikzstyle{every state}=[fill=red!50,draw=none]
      \node[state, scale=0.7] (A)    {$0$};
      \path (A) edge [loop right] node {\footnotesize{(1, 0, $0^\circ$)}} (A)
                edge [loop left] node {\footnotesize{(-1,0,$0^\circ$)}} (A);
    \end{tikzpicture}
    % \caption{Lattice graph A.}
    % \label{fig:lg1}
  \end{minipage}    
  \begin{minipage}[b]{0.5\linewidth}
    \centering
    \begin{tikzpicture}
    \draw[fill=blue] (3,2.92) -- (2.5,2.75) -- (3,2.58) -- (2.875,2.75) -- cycle;
    \draw[fill=blue] (2,2.92) -- (1.5,2.75) -- (2,2.58) -- (1.875,2.75) -- cycle;
    \draw[fill=blue] (1,2.92) -- (0.5,2.75) -- (1,2.58) -- (0.875,2.75) -- cycle;
    \draw[dashed](2.5,2.75) -- (1.5,2.75);
    \draw[dashed](1.5,2.75) -- (0.5,2.75);
    \draw[dashed](0.5,2.75) -- (-0.5,2.75);
    \draw[dashed](2.5,2.75) -- (4,2.75);
  \end{tikzpicture}
  \end{minipage}
  \caption{[left] Lattice graph representing an infinite line. [right] Three robots satisfying the left graph.}
%   \begin{minipage}[b]{0.3\linewidth}
%     \centering
%     \begin{tikzpicture}[->,>=stealth',node distance=5cm]
%       \tikzstyle{every state}=[fill=red!50,draw=none]
%       \node[state, scale=0.7] (A)    {$0$};
%       \path (A) edge [loop above] node {\footnotesize{(0, 1, $45\sqrt{2}^\circ$)}} (A)
%                 edge [loop below] node {\footnotesize{(0, -1, $-45\sqrt{2}^\circ$)}} (A);
%     \end{tikzpicture}
%     \caption{Lattice graph B.}
%     \label{fig:lg2}
% \end{minipage}